\chapter{Ход работы}
\label{ch:practice}

В этой главе по порядку описана реализация интерактивного резюме: каркас
страницы, рендер из конфигурации, редактирование с сохранением, загрузка
фотографии, CSS-анимации, эффект Material Wave, экспорт в PDF и публикация
на сервере. Для каждого этапа приведён код решения и результат.

\section{Каркас страницы}

Файл \texttt{index.html} содержит панель управления \texttt{<header>} с
кнопками <<Скачать PDF>> и <<Сбросить>>, пустой контейнер
\texttt{<main~id="resume-root">} (его заполняет рендер) и подключение трёх
скриптов. Скрипты --- классические, без \texttt{type="module"}, поэтому
страница работает и при открытии напрямую через \texttt{file://}.

\begin{lstlisting}[caption={Ключевая часть \texttt{index.html}},captionpos=b]
<header class="toolbar">
    <span class="toolbar__brand">Конструктор резюме</span>
    <p class="toolbar__hint">Кликните по тексту или фото, чтобы отредактировать</p>
    <div class="toolbar__actions">
        <button id="btn-pdf" class="btn btn--primary ripple-host" type="button">Скачать PDF</button>
        <button id="btn-reset" class="btn btn--ghost ripple-host" type="button">Сбросить</button>
    </div>
</header>

<main id="resume-root" class="resume-wrap"></main>
<input type="file" id="photo-input" accept="image/*" hidden>

<script src="js/config.js"></script>
<script src="js/resume.js"></script>
<script src="js/start.js"></script>
<script>
    document.addEventListener("DOMContentLoaded", function () { start(); });
</script>
\end{lstlisting}

\section{Рендер резюме из конфигурации}

Содержимое резюме задаётся объектом \texttt{RESUME\_DATA} в файле
\texttt{config.js}. Функция \texttt{renderResume} собирает из него
семантическую разметку шаблонными строками: массивы навыков, контактов,
записей опыта и образования преобразуются методом \texttt{map}. Каждому
текстовому узлу присваивается атрибут \texttt{contenteditable="true"} и
уникальный \texttt{data-field}, по которому поле связывается с хранилищем.

\begin{lstlisting}[caption={Рендер записей опыта работы (фрагмент \texttt{resume.js})},captionpos=b]
const expHtml = data.experience.map(function (item, i) {
    return `
        <li class="entry ripple-host">
            <div class="entry__top">
                <h3 class="entry__role editable" contenteditable="true"
                    data-field="exp-${i}-role">${escapeHtml(item.role)}</h3>
                <span class="entry__period editable" contenteditable="true"
                      data-field="exp-${i}-period">${escapeHtml(item.period)}</span>
            </div>
            <p class="entry__company editable" contenteditable="true"
               data-field="exp-${i}-company">${escapeHtml(item.company)}</p>
            <p class="entry__details editable editable--multiline" contenteditable="true"
               data-field="exp-${i}-details">${escapeHtml(item.details)}</p>
        </li>`;
}).join('');
\end{lstlisting}

Все значения пропускаются через функцию \texttt{escapeHtml} --- защита от XSS
при сборке разметки через \texttt{innerHTML}. Резюме свёрстано в две колонки
на \texttt{CSS~Grid}: боковая колонка \texttt{<aside>} (фото, контакты,
навыки, языки) и основная \texttt{<main>} (имя, <<Обо мне>>, опыт,
образование).

% — Скриншот: резюме локально, десктоп —
\begin{figure}[h]
  \centering
  \includegraphics[width=0.85\textwidth]{img/01_resume_desktop}
  \caption{Сайт-резюме, открытый локально (десктопная версия)}
  \label{fig:resume_desktop}
\end{figure}

\section{Редактирование и сохранение}

Функция \texttt{start()} навешивает на каждое поле \texttt{.editable}
обработчики. По событию \texttt{focus} запускается анимация-подсветка, по
\texttt{blur} --- сохранение и анимация <<сохранено>>. Сохранение происходит
не на каждое нажатие клавиши, а один раз при уходе из поля; флаг
\texttt{dirty} отмечает наличие несохранённых правок.

\begin{lstlisting}[caption={Обработчики редактируемого поля (\texttt{start.js})},captionpos=b]
field.addEventListener('focus', function () {
    triggerAnimation(field, 'is-editing');
});
field.addEventListener('input', function () {
    dirty = true;
});
field.addEventListener('blur', function () {
    flushSave();
    triggerAnimation(field, 'is-saved');
});
field.addEventListener('keydown', function (event) {
    if (event.key === 'Enter' && !field.classList.contains('editable--multiline')) {
        event.preventDefault();
        field.blur();
    }
});
\end{lstlisting}

Тексты всех полей собираются в объект и сохраняются в \texttt{localStorage};
при загрузке страницы они восстанавливаются. Чтение и запись обёрнуты в
\texttt{try/catch} --- при недоступном хранилище приложение продолжает
работу.

\begin{lstlisting}[caption={Сбор полей и сохранение (\texttt{resume.js})},captionpos=b]
function collectFields() {
    const fields = {};
    document.querySelectorAll('[data-field]').forEach(function (el) {
        fields[el.getAttribute('data-field')] = el.innerText.trim();
    });
    return fields;
}

function saveResumeData(data) {
    try {
        localStorage.setItem(STORAGE_KEY, JSON.stringify(data));
        return true;
    } catch (e) {
        return false;
    }
}
\end{lstlisting}

Восстановление правок выполняется через \texttt{textContent} (а не
\texttt{innerHTML}) --- сохранённое значение всегда вставляется как обычный
текст, поэтому stored-XSS через \texttt{contenteditable} или вставку из
буфера невозможен.

% — Скриншот: поле в режиме редактирования —
\begin{figure}[h]
  \centering
  \includegraphics[width=0.85\textwidth]{img/02_editing}
  \caption{Редактирование поля резюме с анимацией-подсветкой}
  \label{fig:editing}
\end{figure}

\section{Загрузка фотографии}

Клик по аватару открывает скрытый \texttt{<input type="file">}. Выбранный
файл читается объектом \texttt{FileReader}, затем уменьшается через
\texttt{<canvas>} до 400~пикселей по большей стороне и кодируется в JPEG ---
так фотография укладывается в квоту \texttt{localStorage}.

\begin{lstlisting}[caption={Чтение и сжатие фотографии (\texttt{resume.js})},captionpos=b]
function readPhoto(file, callback) {
    if (!file || file.type.indexOf('image/') !== 0) { return; }
    const reader = new FileReader();
    reader.onload = function () {
        const img = new Image();
        img.onload = function () {
            const maxSide = 400;
            const scale = Math.min(1, maxSide / Math.max(img.width, img.height));
            const canvas = document.createElement('canvas');
            canvas.width = Math.round(img.width * scale);
            canvas.height = Math.round(img.height * scale);
            canvas.getContext('2d').drawImage(img, 0, 0, canvas.width, canvas.height);
            callback(canvas.toDataURL('image/jpeg', 0.85));
        };
        img.src = reader.result;
    };
    reader.readAsDataURL(file);
}
\end{lstlisting}

% — Скриншот: загрузка фотографии —
\begin{figure}[h]
  \centering
  \includegraphics[width=0.85\textwidth]{img/03_photo_upload}
  \caption{Загрузка фотографии в резюме через выбор файла}
  \label{fig:photo}
\end{figure}

\section{CSS-анимации изменений}

Все анимации описаны в CSS правилами \texttt{@keyframes}; JavaScript лишь
переключает классы. Функция \texttt{triggerAnimation} снимает класс, делает
форс-reflow и навешивает класс заново --- так одноразовая анимация
перезапускается при каждом изменении.

\begin{lstlisting}[caption={Анимации входа в правку и сохранения (\texttt{style.css})},captionpos=b]
@keyframes edit-pulse {
    0%   { box-shadow: inset 0 0 0 0 rgba(59, 110, 240, 0); }
    35%  { box-shadow: inset 0 0 0 2px #3b6ef0; }
    100% { box-shadow: inset 0 0 0 0 rgba(59, 110, 240, 0); }
}
@keyframes saved-flash {
    0%   { background-color: rgba(34, 197, 94, 0.45); }
    100% { background-color: transparent; }
}
.is-editing { animation: edit-pulse 420ms ease-out; }
.is-saved   { animation: saved-flash 420ms ease-out; }
\end{lstlisting}

Анимации \texttt{edit-pulse} и \texttt{saved-flash} изменяют разные свойства
(\texttt{box-shadow} и \texttt{background-color}), поэтому корректно
сосуществуют, если классы окажутся на элементе одновременно.

\section{Эффект Material Wave}

<<Волна>> создаётся одним делегированным обработчиком \texttt{pointerdown}
на всём документе. Метод \texttt{closest('.ripple-host')} находит ближайший
элемент-хост; функция \texttt{createRipple} вычисляет диаметр и положение
круга и добавляет \texttt{<span>}, который удаляется по окончании анимации.

\begin{lstlisting}[caption={Создание <<волны>> (\texttt{resume.js})},captionpos=b]
function createRipple(event, host) {
    const rect = host.getBoundingClientRect();
    const size = Math.max(rect.width, rect.height);
    const radius = size / 2;
    const ripple = document.createElement('span');
    ripple.className = 'ripple';
    ripple.style.width = size + 'px';
    ripple.style.height = size + 'px';
    ripple.style.left = (event.clientX - rect.left - radius) + 'px';
    ripple.style.top = (event.clientY - rect.top - radius) + 'px';
    ripple.addEventListener('animationend', function () {
        ripple.remove();
    }, { once: true });
    host.appendChild(ripple);
}
\end{lstlisting}

% — Скриншот: эффект Material Wave —
\begin{figure}[h]
  \centering
  \includegraphics[width=0.85\textwidth]{img/04_ripple}
  \caption{Эффект Material Wave при клике по элементу резюме}
  \label{fig:ripple}
\end{figure}

\section{Экспорт в PDF}

Кнопка <<Скачать PDF>> вызывает \texttt{window.print()}. Перед печатью имя
из резюме записывается в \texttt{document.title} --- оно становится именем
PDF-файла в диалоге <<Сохранить как PDF>>; после печати заголовок
восстанавливается по событию \texttt{afterprint}.

\begin{lstlisting}[caption={Запуск печати (\texttt{start.js})},captionpos=b]
document.getElementById('btn-pdf').addEventListener('click', function () {
    flushSave();
    const previousTitle = document.title;
    document.title = buildPrintFilename(collectFields().name);
    window.addEventListener('afterprint', function restore() {
        document.title = previousTitle;
        window.removeEventListener('afterprint', restore);
    });
    window.print();
});
\end{lstlisting}

Внешний вид PDF задаёт блок \texttt{@media print}: панель управления,
подсказки и <<волны>> скрываются, убираются тени, блоки резюме защищаются от
разрыва между страницами правилом \texttt{break-inside:~avoid}.

% — Скриншот: предпросмотр печати —
\begin{figure}[h]
  \centering
  \includegraphics[width=0.85\textwidth]{img/05_print_preview}
  \caption{Предпросмотр печати: резюме готово к сохранению в PDF}
  \label{fig:print}
\end{figure}

\section{Адаптивность}

Двухколоночная сетка \texttt{CSS~Grid} перестраивается медиавыражениями: при
ширине менее 768~пикселей колонки складываются в одну, при ширине менее
480~пикселей дополнительно уменьшаются заголовки. Свойство
\texttt{grid-template-areas} позволяет менять порядок колонок без правки
разметки.

% — Скриншот: мобильная версия —
\begin{figure}[h]
  \centering
  \includegraphics[width=0.45\textwidth]{img/06_resume_mobile}
  \caption{Адаптивная вёрстка на мобильной ширине (\(360\,\)px)}
  \label{fig:mobile}
\end{figure}

\section{Публикация на сервере}

Готовый сайт публикуется на собственном домашнем сервере за обратным прокси
\textbf{Caddy}. В корне \texttt{lab14/} лежит идемпотентный скрипт
\texttt{deploy.sh} с двумя ключами:

\begin{lstlisting}[caption={Запуск деплоя и сноса сайта},captionpos=b]
./deploy.sh --install     # развернуть сайт
./deploy.sh --uninstall   # снести сайт после показа
\end{lstlisting}

Шаги режима \texttt{--install}:
\begin{enumerate}
  \item Запрашивает домен (по умолчанию ---
        \texttt{resume.server34.netcraze.club}) и путь к исходникам сайта.
  \item Проверяет, что Caddy уже установлен; иначе --- ошибка и выход.
  \item Делает резервные копии \texttt{docker-compose.yml} и
        \texttt{Caddyfile} с меткой времени.
  \item Синхронизирует исходники \texttt{lab14/resume/} в
        \texttt{/opt/stack/resume/site/} командой \texttt{rsync}.
  \item Добавляет (если ещё нет) сервис \texttt{resume} на образе
        \texttt{nginx:alpine} в общий \texttt{docker-compose.yml}.
  \item Дописывает блок поддомена в \texttt{Caddyfile}, запускает контейнер
        и перезагружает Caddy.
\end{enumerate}

Режим \texttt{--uninstall} выполняет обратные действия: останавливает
контейнер, awk-фильтром вырезает сервис из \texttt{docker-compose.yml} и блок
поддомена из \texttt{Caddyfile}, удаляет каталог сайта и перезагружает Caddy.
Первый запрос к новому поддомену запускает выпуск TLS-сертификата
Let's~Encrypt; готовый адрес --- \url{https://resume.server34.netcraze.club}.

% — Скриншот: живой сайт по HTTPS —
\begin{figure}[h]
  \centering
  \includegraphics[width=0.85\textwidth]{img/07_resume_live_https}
  \caption{Развёрнутый сайт по адресу
           \texttt{https://resume.server34.netcraze.club}
           (Let's~Encrypt-сертификат активен)}
  \label{fig:resume_live}
\end{figure}

% — Скриншот: docker compose ps на сервере —
\begin{figure}[h]
  \centering
  \includegraphics[width=0.85\textwidth]{img/08_docker_compose_ps}
  \caption{Сервис \texttt{resume} в общем \texttt{docker-compose.yml}
           домашнего сервера}
  \label{fig:compose_ps}
\end{figure}

\endinput
